\documentclass[10pt,a4paper]{amsart}
\usepackage[margin=19mm]{geometry}
\usepackage{amsmath,amssymb,mathtools}
\usepackage[scr=rsfs,cal=euler]{mathalfa}
\usepackage{xfrac}
\usepackage[unicode,hidelinks,bookmarks=false]{hyperref}
\newcommand{\ie}{i.e.\ }
\newcommand{\cf}{cf.\ }
\newcommand{\resp}{respectively\ }
\newcommand{\Subsection}{\subsection}
\providecommand{\nobreakdash}{\nobreak\hbox{-}}
\setlength{\emergencystretch}{2em}
\multlinegap0pt
\usepackage{bm}
\usepackage{bbm}
\usepackage{dsfont}
\usepackage{xspace}
\usepackage{cancel}
\usepackage{mathtools}
\usepackage{amsthm}
\makeatletter
\@ifundefined{theorem}{\newtheorem{theorem}{Theorem}}{}
\@ifundefined{definition}{\newtheorem{definition}[theorem]{Definition}}{}
\@ifundefined{lemma}{\newtheorem{lemma}[theorem]{Lemma}}{}
\@ifundefined{remark}{\newtheorem{remark}[theorem]{Remark}}{}
\makeatother
\makeatletter
\def\acts{\curvearrowright}
\expandafter\ifx\csname customcor\endcsname\relax
\newenvironment{customcor}[1]
  {\renewcommand\theinnercustomcor{#1}\innercustomcor}
  {\endinnercustomcor}
\fi
\expandafter\ifx\csname customq\endcsname\relax
\newenvironment{customq}[1]
  {\renewcommand\theinnercustomq{#1}\innercustomq}
  {\endinnercustomq}
\fi
\expandafter\ifx\csname customthm\endcsname\relax
\newenvironment{customthm}[1]
  {\renewcommand\theinnercustomthm{#1}\innercustomthm}
  {\endinnercustomthm}
\fi
\makeatother
\title[Minimal operations over permutation groups]{Minimal operations over permutation groups}
\author{Paolo Marimon, Michael Pinsker}
\date{Source: arXiv:2410.22060v3 (CC BY 4.0)}
\begin{document}
\maketitle

\noindent\emph{Source and licence.} Minimal operations over permutation groups. Version arXiv:2410.22060v3 (CC BY 4.0), distributed under the Creative Commons Attribution 4.0 International licence, as recorded in the arXiv licence metadata for this exact version and shown on the abstract page https://arxiv.org/abs/2410.22060v3. Full rights notice: licenses/arxiv-2410.22060-CC-BY-4.0.txt.\par\smallskip

\noindent\textit{Editorial scope.} Counted unit: the Lemma whose source label is \texttt{lem:everythingpalfy} in the article (source0.tex lines 1945--1946, source label \texttt{lem:everythingpalfy}), delivered together with the article's own complete original proof of it (source0.tex lines 1947--1987, one \texttt{proof} environment of 41 lines). No mathematics is written, repaired or re-derived: statement and proof are the article's own text line for line. Required context retained uncounted: def:palfy (lines 1926--1936); rem:stuffinstuff (lines 1938--1943). Every definition, display and label of those lines is retained unchanged. Omitted from this package and not redistributed: 1--1925, 1937--1937, 1944--1944, 1988--2330. Nothing inside a retained excerpt is omitted or altered: every retained body is delivered exactly as the article's lines are written, so no omission receipt of any kind accompanies this package. Citations: the retained excerpts carry no citation command at all, so no bibliography entry of the article is transcribed and no reference is redistributed. Reference adaptation: none was needed and none was made; every reference inside the delivered unit resolves inside it, so no unresolved reference or citation is produced. Printed slips kept literal: no printed word, symbol, hypothesis or number of the counted statement or of its proof was altered. Layout and class: the article's own class is \texttt{article} with packages including geometry, inputenc, mathpazo, hyperref, doi, relsize, fontenc, amssymb, amsmath, xcolor, babel, amsthm, enumitem, tikz; the wrapper is the lane's pinned \texttt{amsart} editorial wrapper at 10pt on A4, which restates the theorem-like environments the retained text needs and adds the article's own macro definitions that the retained text needs (\texttt{\textbackslash acts}), with extra wrapper packages (bm, bbm, dsfont, xspace, cancel, mathtools). The delivered numbering is the unit's own. Assets: no retained span contains a figure environment, a graphics-inclusion command, a PDF-page inclusion, a table or a TikZ picture; no image file is copied into this package, the package declares no asset dependency and no image file is redistributed with it. Licence: version arXiv:2410.22060v3 of this article is distributed under the Creative Commons Attribution 4.0 International licence, as recorded in the arXiv licence metadata for this exact version and shown on the abstract page https://arxiv.org/abs/2410.22060v3. A full rights notice is delivered at licenses/arxiv-2410.22060-CC-BY-4.0.txt. Characters: every retained excerpt is pure ASCII, so the delivered engine, XeLaTeX with its default Latin Modern text font, reports no missing character.
\begin{definition}\label{def:palfy} Let $G\acts B$ with $s$ many orbits. Let $b\in B$ be in a $G$-orbit consisting of more than one element. Let {$\alpha,\beta\in\overline{G}$.} For finite $k\leq s$, we define the $k$-ary orbit\-/semiprojection $f^{(\alpha, \beta)}$ as follows: for $(a_1, \dots, a_k)\in B^k$, 
\begin{equation*}
%\label{eq:palfy}
 f^{(\alpha, \beta)}(a_1, \dots, a_k):=
    \begin{cases}
       \beta b \text{ if } a_1\sim b \text{ and } a_1, \dots, a_k \text{ are all in distinct orbits};\\
        \alpha a_1 \text{ otherwise.}
        \end{cases}
        \end{equation*}
We write $f$ for $f^{(1, 1)}$, where $1$ is the identity in $G$.
\end{definition}

\begin{remark}\label{rem:stuffinstuff} Note that any $f^{(\alpha, \beta)}$ for $\alpha, \beta\in G$ is in $\langle G\cup\{f\}\rangle$, being obtained by
\[\beta f(\beta^{-1}\alpha x_1, \dots, x_k).\] 
Conversely, for any $\alpha, \beta\in G$, $f\in \langle G\cup\{f^{(\alpha, \beta)}\}\rangle$. Also note that for $\gamma\in G$,
\[\gamma f^{(\alpha, \beta)}=f^{(\gamma\alpha, \gamma\beta)}.\]

\end{remark}

\begin{lemma}\label{lem:everythingpalfy} {Let $\alpha, \beta\in G$ and $t\in\langle G \cup \{f^{(\alpha, \beta)}\}\rangle \setminus \langle G\rangle$. Then, $t$ is of the form $f^{(\alpha', \beta')}$ for some $\alpha', \beta'\in G$ up to permuting variables and removing dummy variables.}
\end{lemma}

\begin{proof} We prove this by induction on complexity of $t$, where the base case is trivial. We now carry out the inductive step.
{Consider first the case where $t=\gamma t_1$ for $\gamma\in G$ and $t_1\in\langle G \cup \{f^{(\alpha, \beta)}\}\rangle$ satisfying the induction hypothesis. If $t_1\in\langle G\rangle$ then so is $t$; otherwise,  $t_1=f^{(\alpha',\beta')}$ (up to permuting variables and removing dummy variables), and $t=f^{(\gamma \alpha',\gamma\beta')}$ (with the same caveat) by Remark~\ref{rem:stuffinstuff}.}
%{Firstly, if $t$ is of the form $\gamma\delta$ for some $\gamma, \delta\in G$, then $t\in\langle G\rangle$, and so we are not concerned with this case. Secondly, if $t$ is of the form $\gamma f^{(\alpha', \beta')}$ for some $\gamma, \alpha', \beta'\in G$, from Remark~\ref{rem:stuffinstuff}, we know that $t$ is of the form $f^{(\gamma\alpha', \gamma\beta')}$. }
%For the third case of the inductive step, we need to consider $t$ of the form 
{Secondly, consider $t$ of the form}
$f^{(\alpha, \beta)}(t_1, \dots, t_k)$ where the $t_i\in\langle G \cup \{f^{(\alpha, \beta)}\}\rangle$ {satisfy the induction hypothesis.} %are in $\langle G \cup \{f\}\rangle$. By inductive hypothesis, 
We may suitably permute variables to {assume $t_1$} is %obtain $t'$ 
either of the form %$f^{(\alpha, \beta)}(s_1, \dots, s_k)$, where $t_1$ is either of the form 
$\gamma x_1$ or of the form $f^{(\gamma, \delta)}(x_1, \dots, x_k)$ for some $\gamma, \delta\in G$. Also, %by inductive hypothesis 
 each of the $t_i$ for $i>1$ is of the form $\gamma x_{\tau(i)}$ for $\tau(i)\in\mathbb{N}$ and $\gamma_i\in G$, or of the form $f^{(\gamma_i, \delta_i)}(x_{\sigma^i(1)}, \dots, x_{\sigma^i(k)})$, where $\sigma^i$ is some injection $\{1, \dots, k\}\mapsto\mathbb{N}$, and $\gamma_i, \delta_i\in G$. For $i>1$, let $l(t_i)$ be $\tau(i)$ in the first case and $\sigma^i(1)$ in the second case. The identities in the following claim then tell us that {$t$} is either in $\langle G\rangle$ or again of the form $f^{(\alpha', \beta')}$ for $\alpha', \beta'\in G$, completing the proof.\\

\textbf{Claim:} we have the following identities:
\begin{itemize}
    \item If $l(t_i)=i$ for all $2\leq i\leq k$,
    \begin{equation}\label{eq:palfy1}
        f^{(\alpha, \beta)}(\gamma x_1, t_2, \dots, t_k)\approx f^{(\alpha \gamma, \beta)};
    \end{equation}
    \item If $l(t_i)=1$ or $l(t_i)=l(t_j)$ for some $2\leq i<j\leq k$, 
    \begin{equation}\label{eq:palfy2}
        f^{(\alpha, \beta)}(\gamma x_1, t_2, \dots, t_k)\approx \alpha\gamma x_1;
    \end{equation}
    \item If $l(t_i)=i$ for all $2\leq i\leq k$,
    \begin{equation}\label{eq:palfy3}
        f^{(\alpha, \beta)}(f^{(\gamma, \delta)}, t_2, \dots, t_k)\approx f^{(\alpha \gamma, \beta)};
    \end{equation}
     \item If $l(t_i)=1$ or $l(t_i)=l(t_j)$ for some $2\leq i<j\leq k$, 
    \begin{equation}\label{eq:palfy4}
        f^{(\alpha, \beta)}(f^{(\gamma, \delta)}, t_2, \dots, t_k)\approx \alpha f^{(\gamma, \delta)}.
    \end{equation}
\end{itemize}
\begin{proof}[Proof of claim] The identities follow directly from Definition~\ref{def:palfy}. For example, for (\ref{eq:palfy1}), given $(a_1, \dots, a_k)$, we have that
\[f^{(\alpha, \beta)}(\gamma x_1, t_2, \dots, t_k)(a_1, \dots, a_k)=f^{(\alpha, \beta)}(\gamma a_1, a_2', \dots, a_k'),\]
where $a_i'\sim a_i$. Hence, if $(a_1, \dots, a_k)$ were all in distinct orbits and $a_1\sim b$, we get that 
\[f^{(\alpha, \beta)}(\gamma a_1, a_2', \dots, a_k')=\beta b.\]
Otherwise, we have
\[f^{(\alpha, \beta)}(\gamma a_1, a_2', \dots, a_k')=\alpha\gamma a_1.\]
This yields the desired identity in (\ref{eq:palfy1}). For identity (\ref{eq:palfy2}), we get that 
\[f^{(\alpha, \beta)}(\gamma x_1, t_2, \dots, t_k)(a_1, \dots, a_k)=f^{(\alpha, \beta)}(\gamma a_1, a_{l(t_2)}', \dots, a_{l(t_k)}'),\]
where $a_{l(t_k)}'\sim a_{l(t_k)}$. In particular, since either $l(t_i)=1$ or $l(t_i)=l(t_j)$ for $2\leq i<j\leq k$, we are never in the first condition of Definition~\ref{def:palfy}, and so the identity (\ref{eq:palfy2}) holds. The final two identities (\ref{eq:palfy3}) and (\ref{eq:palfy4}) follow by the same reasoning as the previous two.
\end{proof}
\end{proof}

\end{document}
